\begin{table}[H]
\centering
\caption{Remuneración laboral de ocupados de 25 a 29 años con educación universitaria o superior, 2010 y 2025}
\label{tab:remuneracion_recien_graduados}
\footnotesize
\begin{tabular}{@{}p{6.6cm}rrr@{}}
\toprule
Indicador & 2010 & 2025 & Crec. anual \\
\midrule
Ocupados (millones de personas) & 0,73 & 1,27 & 3,73\% \\
Remuneración mensual por trabajador (millones de pesos de 2025) & 2,33 & 2,40 & 0,18\% \\
Remuneración por hora trabajada (miles de pesos de 2025) & 11,7 & 12,7 & 0,51\% \\
\bottomrule
\end{tabular}
\caption*{\footnotesize Nota: la GEIH no identifica si la persona acaba de graduarse ni el año de graduación. Este cuadro aproxima la remuneración de recién graduados con personas ocupadas de 25 a 29 años que reportan educación universitaria o superior. Por lo tanto, debe leerse como una aproximación a jóvenes con educación superior, no como graduados observados directamente. El crecimiento es anualizado para 2010--2025. La serie excluye 2020. Fuente: cálculos propios con GEIH del DANE.}
\end{table}
